%% ============================================================
%% §6 — Discussion
%% Merged P1+P6 Journal Paper: SiliQun — A Physics-Grounded
%% Simulation Platform and Framework for the PGIRS Data Plane
%%
%% Sources:
%%   §6.1 F≥0.99 Significance     → P1 §5.1 (adapted, reframed as SiliQun validation)
%%   §6.2 Positioning vs Related Work → P1 §5.2 + P6 §2.4 Oxford framing
%%   §6.3 PGIRS Generalisability   → NEW (beyond SiMOS applicability)
%%   §6.4 a hierarchical DRL framework Scalability    → NEW (forward ref, placeholder for results)
%%   §6.5 Limitations + Future Work → P1 §5.3 (augmented with sim-to-hardware gap)
%% ============================================================

\section{Discussion}
\label{sec:discussion}

%% ----------------------------------------------------------
\subsection{Significance of the F \texorpdfstring{$\geq$}{>=} 0.99 Result as a Simulator Validation}
\label{subsec:significance}
%% ----------------------------------------------------------

The achievement of $F_\text{eval} = 0.9930$ by the DRL control agent on SiliQun is significant not only as a DRL result but as a validation of the simulation platform itself. The fault-tolerance threshold $F \geq 0.99$ is not an arbitrary benchmark: it corresponds to the fidelity above which the surface code can suppress logical error rates below the physical error rate~\cite{Fowler2012}, making it the practical boundary between NISQ-era and fault-tolerant quantum computing. Demonstrating that a DRL agent trained exclusively on SiliQun can cross this threshold under first-principles SiMOS Lindblad noise establishes two things simultaneously.

The first is that SiliQun's physics engine is accurate enough to pose a genuinely hard optimisation problem. The Fix3 baseline — a competent DRL agent with manually tuned reward weights — fails to cross $F = 0.99$ after 300,000 training steps under the same SiMOS noise model. If the noise model were too weak or too easily circumvented, Fix3 would have succeeded. The fact that it does not confirms that SiliQun's Lindblad dynamics are physically realistic enough to expose the true difficulty of 4-qubit gate synthesis on SiMOS hardware.

The second is that SiliQun's Gymnasium interface is well-formed enough to support the full the DRL control agent algorithm stack — SAC, BRFD, HMRL, DEHB, expert warm-start, and curriculum learning — without any environment-specific modifications. This is a non-trivial property: several of the seven DRL families benchmarked in \S\ref{subsec:cross_algo} fail to converge or converge slowly, not because of environment defects but because of fundamental mismatches between the algorithm and the task structure. The fact that the most sophisticated algorithm (SAC-HMRL) achieves $F = 0.993$ while the simplest (DQN) achieves only $F = 0.847$ confirms that the environment provides a rich enough learning signal to discriminate between algorithm families — a prerequisite for a useful benchmark environment.

The rapidity of convergence — 194 seconds on a single A100 GPU — is attributable to three design choices that are themselves part of the PGIRS methodology (\S\ref{subsec:pgirs}): expert warm-start providing accurate Q-value initialisation, BRFD discovering well-calibrated reward weights for the 4-qubit landscape, and the curriculum patches preventing catastrophic forgetting events that caused earlier runs to regress after reaching high fidelity. The 93$\times$ GPU speedup over CPU simulation (\S\ref{subsec:gymnasium}) is a necessary prerequisite for this convergence speed: at CPU speeds, the 20,000 training steps required to reach $F = 0.99$ would take approximately 5 hours rather than 194 seconds, making iterative algorithm development impractical.

%% ----------------------------------------------------------
\subsection{Positioning Within the Quantum Simulation Landscape}
\label{subsec:positioning}
%% ----------------------------------------------------------

SiliQun occupies a specific and previously unoccupied niche in the quantum simulation landscape. Three positioning dimensions clarify this niche.

\textbf{Against general-purpose quantum simulators (QuTiP, Qiskit Aer, PennyLane).} These frameworks provide accurate quantum dynamics simulation but do not expose a Gymnasium-compatible interface. Converting them to RL training environments requires substantial engineering effort — defining observation and action spaces, implementing reward functions, managing episode boundaries, and integrating with SB3 or other RL libraries. SiliQun performs this integration once, correctly, and releases it as a reusable open-source package. The PGIRS methodology (\S\ref{subsec:pgirs}) codifies the engineering decisions made during this integration so that future researchers extending SiliQun to new hardware platforms can follow the same five-phase discipline.

\textbf{Against abstract RL environments for quantum circuit synthesis (qgym, Qiskit Gym, QuantumCircuitEnv).} These environments provide Gymnasium-compatible interfaces but use abstract depolarising noise models that do not capture the amplitude damping and dephasing channels present in SiMOS devices. As demonstrated in \S\ref{subsec:cross_algo}, the choice of noise model has a substantial impact on achievable fidelity: algorithms that solve the task trivially under abstract depolarising noise (DQN, PPO) fail to reach the fault-tolerance threshold under SiMOS Lindblad noise. SiliQun is the first environment to combine a Gymnasium-compatible interface with a first-principles SiMOS Lindblad noise model.

\textbf{Against the Oxford DPhil spin qubit work (Orbell 2024~\cite{Orbell2024}).} As detailed in \S\ref{subsec:oxford_comparison}, Orbell's work addresses the device tuning and calibration layer of the quantum computing stack, while SiliQun addresses the gate synthesis layer. The two works are non-overlapping in problem domain, noise modelling approach, and software availability. The natural integration path is for Orbell's Bayesian optimisation pipeline to calibrate SiliQun's hardware profile from experimental data, replacing the current manually chosen conservative parameters with device-specific values. This would realise the full simulation-to-synthesis pipeline from physical hardware characterisation through gate synthesis.

%% ----------------------------------------------------------
\subsection{PGIRS Generalisability Beyond SiMOS}
\label{subsec:generalisability}
%% ----------------------------------------------------------

The PGIRS methodology (\S\ref{subsec:pgirs}) was developed in the context of SiMOS spin qubits, but its five phases are hardware-agnostic by design. Each phase operates on the device profile $\mathcal{P}$ as an abstract parameter set; replacing the SiMOS profile with a different hardware profile instantiates a different simulation environment without modifying the framework architecture.

Three hardware platforms are natural candidates for PGIRS extension. \textbf{Superconducting transmon qubits} (IBM, Google, Rigetti) are the most widely deployed NISQ hardware and have well-characterised noise models~\cite{Krantz2019}. The dominant noise channels — $T_1$ relaxation, $T_2$ dephasing, and gate crosstalk — map directly to the Lindblad operators already implemented in SiliQun's physics engine; extending SiliQun to transmon qubits would require updating the hardware profile $\mathcal{P}$ and the gate set, but not the Lindblad solver or the Gymnasium interface. \textbf{Trapped-ion qubits} (IonQ, Honeywell) have longer coherence times but slower gate speeds and all-to-all connectivity; the cross-backend validation in \S\ref{subsec:cross_backend} already demonstrates that the DRL agent-synthesised circuits transfer to IonQ Aria-1 with $\mathcal{F}_\mathrm{XEB} \geq 0.954$, suggesting that a trapped-ion hardware profile for SiliQun would be a productive direction. \textbf{Neutral atom qubits} (QuEra, Pasqal) are an emerging platform with native multi-qubit gates and reconfigurable connectivity; extending SiliQun to neutral atoms would require adding Rydberg blockade interactions to the Lindblad solver, which is a well-defined extension of the existing framework.

The PGIRS methodology's Phase 5 (empirical validation) provides a systematic protocol for validating any new hardware profile: the three analytical benchmarks (Rabi oscillations, $T_1$ decay, Bell state fidelity) are hardware-independent and can be applied to any Lindblad-based simulation. This provides a reproducible quality gate for hardware profile development that does not depend on access to physical hardware.

%% ----------------------------------------------------------
\subsection{Scalability: From 4-Qubit Gates to Larger Systems}
\label{subsec:scalability}
%% ----------------------------------------------------------

The experiments in this paper focus on 4-qubit GHZ state preparation under SiMOS Lindblad noise, which is the natural entry point for a first-principles silicon spin qubit simulator. Two observations from the experimental results inform the scalability outlook.

\textbf{MPS scaling is near-linear.}
The throughput benchmark in \S\ref{subsec:setup} demonstrates that SiliQun's MPS-based simulation scales near-linearly with qubit count, in contrast to the exponential scaling of statevector simulators. At 8 qubits, SiliQun achieves $6\times$ to $41\times$ speedup over Qiskit Aer and QuTiP, and this advantage grows with qubit count. The MPS bond dimension $\chi$ is the primary control parameter: for the 4-qubit GHZ synthesis task, $\chi = 2$ is sufficient throughout training, and the Lindblad noise overhead remains below 39\% relative to the noiseless MPS baseline. These scaling properties make SiliQun a practical simulation platform for 8- to 16-qubit gate synthesis experiments, which represent the next natural step beyond the 4-qubit results reported here.

\textbf{Policy compositionality enables modular scaling.}
The cross-backend validation in \S\ref{subsec:cross_backend} demonstrates that DRL policies trained on SiliQun transfer to real quantum hardware without retraining. The GAA device extension (\S\ref{subsec:gaa}) further demonstrates that policies trained on the SiMOS profile generalise to a structurally different device with $F = 0.994 \pm 0.012$, confirming that SiliQun-trained policies are not overfit to a single hardware configuration. Extending SiliQun to 8-qubit and 16-qubit training environments is the primary engineering prerequisite for modular multi-qubit gate synthesis, and is identified as a near-term development priority.

%% ----------------------------------------------------------
\subsection{Limitations and Future Work}
\label{subsec:limitations}
%% ----------------------------------------------------------

Four limitations of the current work motivate future research directions.

\textbf{Simulation-to-hardware gap.} The cross-backend validation in \S\ref{subsec:cross_backend} demonstrates that the DRL agent-synthesised circuits achieve $\mathcal{F}_\mathrm{XEB} \geq 0.954$ on IonQ Aria-1, but this is a transferability result, not a hardware calibration result. The SiMOS noise parameters used in SiliQun ($T_1 = 1$~ms, $T_2 = 0.5$~ms, $p_2 = 1\%$) are calibrated to the parameter regime of SiMOS devices as characterised in the literature~\cite{Xue2022,Noiri2022,Philips2022,Steinacker2023}, but have not been validated against a specific physical SiMOS device through standard benchmarking circuits. Closing this gap requires access to a physical SiMOS device — currently available only at a small number of academic and industrial laboratories worldwide — and is identified as the highest-priority future work item.

\textbf{Statistical robustness.} The primary result ($F = 0.9930$) is reported as mean $\pm$ std over five seeds, which is sufficient for a proof-of-concept demonstration but falls short of the statistical rigour expected for a definitive benchmark result. A systematic multi-seed campaign with 20--50 seeds and formal hypothesis testing (Wilcoxon signed-rank test against the Fix3 baseline) is needed to establish the result's statistical significance at the $p < 0.01$ level.

\textbf{Scaling beyond 4 qubits.} The direct training approach used in this paper (single SiliQun instance, 4-qubit GHZ target) does not scale to 8 or 16 qubits without modifications to the state representation and reward function. Extending SiliQun to support 8-qubit and 16-qubit training environments --- with corresponding increases in the density matrix state space and the Lindblad solver's computational cost --- is a necessary prerequisite for multi-qubit gate synthesis beyond the 4-qubit regime.

\textbf{Gate set completeness.} The current SiliQun gate set (11 tokens: 6 single-qubit rotations, 1 two-qubit entangling gate, 4 identity/barrier tokens) is sufficient for GHZ state preparation but does not cover the full universal gate set required for arbitrary quantum computation. Extending the gate set to include $T$ gates, $S$ gates, and controlled-$Z$ gates — and validating the extended gate set against the analytical benchmarks in \S\ref{subsec:validation} — is a straightforward extension of the current framework.

%% ----------------------------------------------------------
%% Bibliography entries for this section (to be merged into main .bib)
%% ----------------------------------------------------------
%%
%% @article{Fowler2012,
%%   author  = {Fowler, Austin G. and Martinis, John M. and others},
%%   title   = {Surface codes: Towards practical large-scale quantum computation},
%%   journal = {Physical Review A},
%%   volume  = {86},
%%   pages   = {032324},
%%   year    = {2012},
%%   doi     = {10.1103/PhysRevA.86.032324}
%% }
%%
%% @article{Krantz2019,
%%   author  = {Krantz, Philip and Kjaergaard, Morten and Yan, Fei and Orlando, Terry P. and Gustavsson, Simon and Oliver, William D.},
%%   title   = {A quantum engineer's guide to superconducting qubits},
%%   journal = {Applied Physics Reviews},
%%   volume  = {6},
%%   pages   = {021318},
%%   year    = {2019},
%%   doi     = {10.1063/1.5089550}
%% }
